\documentclass[a4paper,11pt]{article}
\usepackage{listings}
\usepackage{CJKutf8}
\usepackage{minted}
\usepackage{picinpar,graphicx}
\usepackage{CJK}
\usepackage[cache=false, outputdir=build]{minted}
\author{langman}
\title{ACM-template}

\begin{document}
\maketitle
\vspace{3cm}
\begin{CJK}{UTF8}{gbsn}
    \begin{figure}[!htb]
      \begin{center}
        \includegraphics[width=0.80\linewidth]{../scoure/tupian.jpg}
        \caption{stay hungry stay foolish}
      \end{center}
    \end{figure}
\newpage
\tableofcontents
\newpage
\section{头文件}
\inputminted{c++}{../scoure/head.cpp}
\newpage
\section{图论}
\subsection{二分图}
判断是否二分图
\inputminted{c++}{../scoure/Graph_theory/pan2fen.cpp}
求出最大匹配数
\inputminted{c++}{../scoure/Graph_theory/2fentu.cpp}
\subsection{并查集}
\inputminted{c++}{../scoure/Graph_theory/bingchick.cpp}
\newpage
\subsection{最短路}
两种算法 但是要注意dijkstra无法处理负边的情况
\subsubsection{dijkstra}
需要注意的在于 可以更优化 我没写了 而且需要注意重边的情况
\inputminted{c++}{../scoure/Graph_theory/dijkstra.cpp}
\subsubsection{spfa}
需要注意的是怎么建边 双向边?
\inputminted{c++}{../scoure/Graph_theory/spfa.cpp}
\subsubsection{Flody}
这个就不写了,一个小dp
\subsection{最小生成树}
这是个什么玩意呢 图里面是吧,找到n-1条边使得生成一颗树,然后他的边权之和最小
\inputminted{c++}{../scoure/Graph_theory/prime.cpp}
\subsection{最大流}
\subsubsection{Dinic}
板子先存着,坑定用的着
\inputminted{c++}{../scoure/Graph_theory/dinic.cpp}
\newpage
\section{常用数据结构}
\subsection{STL}
\inputminted{c++}{../scoure/date/duilie.cpp}
\subsection{树状数组}
适用于的方面是，单点更新，多次查询。
\inputminted{c++}{../scoure/date/treevector.cpp}
\subsection{前缀和}
适用于多次区间更新，一次查询，代码就不搞上去了。
\subsection{线段树}
现在的我对于这方面还不够熟悉，先留个板子
\inputminted{c++}{../scoure/date/treexian.cpp}
\subsection{高精度}
我个人习惯使用Java。
\inputminted{java}{../scoure/date/bigjing.java}
\section{数学方面}
\subsection{三个特别的数}
\subsubsection{Fib 数列}
$$f(x) = f(x-1)+f(x-2)$$
$$f(0) = 0,f(1) = 1$$
\subsubsection{卡特兰 数}
$$\sum_{i=1}^n f_i*f_{n-i}=f_n$$
$$h(n) = C_{2n}^n - C_{2n-1}^{n}$$
注意它这个数字来自于什么情况。
\subsubsection{斯大林公式}
$$ \sqrt{2*PI*n} * (\frac{n}{e})^n = n!$$
\subsection{数论}
第一个自然是最基础的欧几里得算法，欧几里得算法的用处有很多，求最大公倍数，解方程，很多。在后面的过程会把一些常见的板子列出来，一般来说这些板子
都已经经过验证，但是不好说对吧。简单题我们可以通过
一些模板直接得出答案，但是怎么说，这些对于难题
估计只能算工具，重要的是如何转换。
\subsubsection{素数}
素数筛法
\inputminted{c++}{../scoure/math/shai.cpp}
\subsubsection{梅森素数}
一个知识点吧，m是一个正整数，且$2^m-1$为素数，那么m一定为素数。\\
如果m是一个素数，$M_p = 2^p-1$是梅森数\\
如果p是一个素数，并且$M_p = 2^p-1$也是素数，那么称$M_p$为梅森素数\\
对梅森素数的判定是一个算法：\\
Lucas-Lehmer:
$r_k \equiv r_{k-1} -2$\quad mod$M_p$\quad $r_1 = 4$
当且仅当 $r_{p-1} \equiv 0$\quad mod$M_p$
\subsubsection{欧几里得}
\inputminted{c++}{../scoure/math/GCD.cpp}
然后是基于这个定理得出的一个定理，中国剩余定理
\inputminted{c++}{../scoure/math/chineseshengyu.cpp}
\subsubsection{乘法逆元}
思想是通过扩展欧几里得来得出，如果缘分到了，那么
还能用费马小定理来解
\inputminted{c++}{../scoure/math/niyuan.cpp}
\subsubsection{欧拉函数}
p为素数时$\phi(p) = p-1 $\\
a与n互质的时候 $a^{\phi(n)}\equiv 1$ mod $ n$\\
m,n互质$\phi(mn) = \phi(m) * \phi(n)$\\

这个东西好呀，他求的是比n小的，并且和n互质的数的个数
\inputminted{c++}{../scoure/math/oula.cpp}
更多的来说我觉得这个东西是一个工具，他对解一些题有很重要的作用，起到一个工具的作用
我目前学的比较浅，对他的优化作用没有很深的了解。
\subsubsection{莫比乌斯函数}
$$ F_n = \sum_{i=1}^n f_i $$
$$ f_i = \sum_{d|n} u(d)*f(\frac{d}{n}) $$
和欧拉函数一样很重要的一个函数他的定义我就不说了，毕竟我latex学的还不好，公式的
插入对我来说用处不大。
\inputminted{c++}{../scoure/math/mobius.cpp}
\newpage
\subsection{博弈}
\subsubsection{主要的解题思想}
官方说的是通过必败点和必胜点来判定
先通过必败点来推，直接来看必胜点，把问题抽象成图 把状态抽象成点，必败点就是先手必败点，然后通过必败点能走到的搞成必胜点，如过有一个状态没有走过 而且他后面的路都是必胜点那么他就是必败点。感觉就像dp一样，记忆化搜索。
当然题目不可能出的那么简单的。
不过根据雄爷定理，万事不离期宗，掌握基本，扩展自己去发掘。
\subsubsection{题型}
巴什博弈\\
这个是最简单的博弈，就是一堆东西，每个人自己能拿1-n件，谁最后一个拿完谁赢，这个是最简单的，不记录。
\\威佐夫博弈\\
有两堆各若干个物品，两个人轮流从某一堆或同时从两堆中取同样多的物品，规定每次至少取一个，多者不限，最后取光者得胜。
这个的解题思路在于通过前面的那个np问题来解决，用局势来思考这些问题，前几个局势在于(0,0),(1,2),(3,5),(4,7).....然后一些大佬就总结出了一些牛逼的结论$( a_k,b_k),a_k=\frac{k*(\sqrt{5}+1)}{2} , b_k=a_k+k$人才。
\\Fibonacci\\
有一堆个数为n的石子，游戏双方轮流取石子，满足：
\\（1）先手不能在第一次把所有的石子取完；
\\（2）之后每次可以取的石子数介于1到对手刚取的石子数的2倍之间（包含1和对手刚取的石子数的2倍）。 约定取走最后一个石子的人为赢家。
结论是 当n为Fibonacci数时，先手必败
\\尼姆博弈\\
有三堆各若干个物品，两个人轮流从某一堆取任意多的物品，规定每次至少取一个，多者不限，最后取光者得胜。
这个博弈有点意思 他的必败点的局势在于$(a,b,c) a \land {b \land c} = 0$
\subsubsection{SG函数}
这个在看之前感觉很高级但是啊，好像也就是一个dp的过程，通过一个必败点，看成起点然后，那个方法看成通向下一个起点的路，然后找所有能直接到这个必败点的必胜点。好像也就那么回事。好像能解决的都是小数字题这是一个板子，f里面存的是方法，多堆问题可以转化成异或来解决。
\inputminted{c++}{../scoure/math/boyi.cpp}
\subsubsection{解题策略}
* 1  ：相信自己的第一感觉\\
* 2　：博弈都会和一些特别的数搭边 ， 所以第一件事坑定是分析局势然后找找看是不是有特别的意义，像什么 卡特兰数，ｆｉｂ数列 ，幂次方，异或的值是否为０；\\
* 3  : 不挂怎么说，记得打表。\\
\subsection{组合数学}
\subsubsection{求组合数}
第一个是求组合数,方法很多不去列举，注意的是一般来说，组合数都是需要去
模一个数，所以他的分母在计算的时候是需要去求逆元的
\inputminted{c++}{../scoure/math/zuhe.cpp}
\subsubsection{polay定理}
设G={p1，p2，…，pt}是X={a1,a2,…,an}上一个置换群，用m种颜色对X中的元素进行涂色，那么不同的涂色方案数为
$$\frac{1}{G}\sum_{k=1}^t m^{Cyc(p_k)}$$
$Cyc(p_k)$是置换$p_k$的循环节个数
\subsubsection{Pell方程}
$$x^2-dy^2=1$$
当d不为平方数时，有无穷多的解。接下来是通项的推导
$$ x_{n+1} = x_1*x_n+d*y_1*y_n,y_{n+1} = x_1*y_n + y_1*x_n $$
暴力求$x_1,y_1$然后矩阵快速幂
\subsubsection{lucas定理}
当组合数的基数过大的时候进行这些操作但是注意，我们的操作也是要求那个模数为素数，且
模数要小的情况下，素数的情况我们可以用扩展lucas定理来解决。一个工具，一个数论上的分支。
\inputminted{c++}{../scoure/math/lucas.cpp}
\subsection{线代}
\subsubsection{fft优化多项式乘法法}
这个是一个工具，他的作用是加速多项式得乘法。这个点得难处还是在于多项式的构建，其实也不是
非常复杂得东西，当然他的原理，不敢去碰。关键词：多项式乘法，并且复杂度为1e5左右。
\inputminted{c++}{../scoure/math/fft.cpp}
\subsubsection{矩阵快速幂}
这类方法，很多是用在递推关系式的时候，像什么fib数列什么的。
\inputminted{c++}{../scoure/math/ju_quick.cpp}
\subsubsection{矩阵方面知识}
就是用高斯消元法去解决一些问题，像什么秩和方阵的值。
\subsection{计算几何}
\subsubsection{几何基本知识}
\emph{矢量}\\
矢量的乘积有很多的作用，注意定义。
适用点在于：1:面积 2：位置\\
\emph{跨立实验与判断两线段是否相交}\\
线段$P_1P_2,Q_1Q_2$,相交的条件为\\
$P_1Q_1$x$P_1P_2$  *  $P_1P_2$x$P_1Q_2>=0$\\
$Q_1P_1$x$Q_1Q_2$  *  $Q_1Q_2$x$Q_1P_2>=0$\\
\emph{pick定理}\\
线段上的是整数点的数的个数 求gcd;\\
PICK定理 设以整数点为顶点的多边形的面积为S，多边形内部的整数点数为N，多边形边界上的整数点数为L，则 S=L/2 + N-1

\subsubsection{判断点是否在多边形中}
\inputminted{c++}{../scoure/math/jihe.cpp}
为1的时候，则在内部。2，应该是边上。
\subsubsection{凸包问题}
\inputminted{c++}{../scoure/math/tubao.cpp}
\section{状态转移 dp}
dp的定义：
1：记忆化搜索；2：状态转移
所以我们的解决方案总是跟着这个来走，从定义出发。
难点集中于两个方面，状态式的确定和状态转移方程的确定
\subsection{背包}
背包的问题主要以下几种：
01背包，部分背包，完全背包;[相对来说比较简单] ，分组背包[个人感觉较难]
背包难在如何， 确定维数，确定背包的容量是什么以及背包的价值是什么，还有背包的dp关系转移式。
\subsection{一些常见的dp}
\subsubsection{LIS 最长上升子序列}
\inputminted{c++}{../scoure/dp/longerxuelie.cpp}
\subsection{树形dp}
关键点在于找状态点间的关系，他一般只有三个关系，父亲节点
，儿子节点，还有兄弟节点，去找他们之间的关系，所以一般是两遍dfs
找父亲与儿子的关系，找儿子与父亲的关系。
\subsection{数位dp}
这个dp的精髓在于记忆化搜索，也就是在最高位不是被限定的情况下进行记录，这样的话省掉很多多余的步骤。
所有的出发点都处于这个目的。
\subsection{状压dp}
这个dp的精髓在于状态转移，不过能压缩的情况也是很限定的。
像什么每个点的状态在于都是能用两个状态来描述，且这些点不多，但是组合的方式很多。
一些状压dp经常用的上的公式。
\inputminted{c++}{../scoure/dp/zhuangtai.cpp}
\subsection{the end}
\end{CJK}
\end{document}
